\chapter{Analiză și fundamentare teoretică}
\label{ch:analysis}
\pagestyle{fancy}

Capitolul de față stabilește cadrul teoretic pe care se sprijină soluția propusă. Sunt explicate, în ordine, problema abordată, elementele de protocol din care provin caracteristicile analizate, modelul de amenințare care delimitează ce poate face un atacator, condiția de consistență pe care trebuie să o respecte orice variantă perturbată, procedurile propuse, argumentele care justifică alegerea celor trei modele comparate, metricile potrivite unui cadru puternic dezechilibrat și structura logică a soluției. Accentul este pus pe motivarea alegerilor: fiecare decizie este prezentată împreună cu alternativa respinsă și cu motivul respingerii.

\section{Problema abordată}\label{sec:formalizare}

\subsection{Ce vede un sistem de detecție}\label{subsec:clasificare}

Un sistem de detecție a intruziunilor care lucrează la nivel de flux nu analizează pachete individuale. Senzorul de rețea urmărește toate pachetele schimbate în cadrul unei conexiuni și produce, la încheierea acesteia, o singură înregistrare care rezumă statistic ceea ce s-a întâmplat: câte pachete și câți octeți a trimis fiecare parte, cât a durat conexiunea, ce valori au avut anumite câmpuri din antete, cât de mari au fost intervalele dintre pachete și cât de intensă a fost activitatea recentă din rețea în jurul perechii de calculatoare implicate.

Această înregistrare, formată în cazul de față din patruzeci și două de mărimi, este singurul lucru pe care îl primește modelul. O parte dintre mărimi sunt numerice, precum volumul de octeți transmis, iar altele sunt categorice, precum protocolul folosit sau starea în care s-a încheiat conexiunea.

Sarcina modelului este să atribuie fiecărei înregistrări una dintre zece etichete posibile: fie traficul este considerat legitim, fie este încadrat în una dintre cele nouă categorii de atac. Modelul răspunde, prin urmare, nu doar la întrebarea dacă traficul este malițios, ci și la întrebarea ce fel de atac reprezintă.

Alegerea unei formulări cu zece categorii, în locul uneia binare care ar distinge doar între trafic legitim și trafic malițios, este esențială pentru obiectivele lucrării. Ea permite distincția dintre două situații pe care o formulare binară le-ar confunda: un flux care nu mai este deloc recunoscut ca malițios și un flux care este în continuare semnalat, dar încadrat în categoria greșită. Consecințele practice ale celor două sunt complet diferite. În primul caz atacul trece nedetectat și nimeni nu află despre el. În al doilea caz, operatorul primește totuși o alertă, chiar dacă aceasta poartă o etichetă eronată.

\subsection{Distanța dintre atacator și model}\label{subsec:distanta}

Observația care structurează întreaga lucrare privește diferența dintre nivelul la care acționează atacatorul și nivelul la care decide modelul.

Atacatorul manipulează pachete. Le poate mări adăugând conținut, le poate întârzia, le poate schimba anumite câmpuri din antet. Modelul, în schimb, nu vede niciodată pachetele: vede doar rezumatul statistic calculat de senzor. Între cele două se interpune procedura de extragere a caracteristicilor.

Consecința acestei interpuneri este că o singură acțiune a atacatorului se traduce, în rezumatul pe care îl primește modelul, prin modificarea simultană a mai multor mărimi. Dacă atacatorul adaugă octeți de umplutură, nu crește doar volumul transmis. Se schimbă automat și dimensiunea medie a unui pachet, pentru că aceasta se obține împărțind volumul la numărul de pachete, și se schimbă și debitul, pentru că acesta se obține împărțind volumul la durată. Practic nicio acțiune realizabilă asupra traficului nu modifică o singură mărime din rezumat.

Această propagare, tratată pe larg în secțiunea~\ref{sec:consistenta}, constituie principala dificultate teoretică a lucrării. Ea explică totodată de ce nu este suficient ca perturbările să fie construite modificând direct valorile din rezumat: un asemenea rezultat nu ar corespunde niciunui trafic pe care un senzor l-ar putea înregistra vreodată.

\subsection{Ce înseamnă evaziunea}\label{subsec:evaziune}

Pornind de la observațiile anterioare, obiectivul atacatorului se poate enunța precis.

Atacatorul dispune de un flux malițios care, în forma sa obișnuită, este detectat. El caută o modificare a traficului cu două proprietăți simultane. Prima: fluxul rezultat trebuie să fie clasificat de model drept trafic legitim. A doua: fluxul rezultat trebuie să fie unul pe care atacatorul îl poate produce efectiv, fără ca atacul să înceteze să funcționeze.

Întreaga dificultate a problemei stă în a defini corect a doua condiție. Fără ea, problema ar fi banală și lipsită de orice interes practic: dacă valorile din rezumatul statistic ar putea fi schimbate liber, orice flux ar putea fi transformat într-unul clasificat drept legitim, prin simpla înlocuire a valorilor cu cele tipice traficului normal. Rezultatul unei asemenea operații nu ar corespunde însă niciunui trafic real, iar concluzia obținută nu ar spune nimic despre securitatea sistemului evaluat.

Mulțimea variantelor pe care atacatorul le poate produce efectiv va fi numită în continuare \emph{mulțimea perturbărilor admisibile}. Construcția ei explicită constituie contribuția metodologică centrală a lucrării și face obiectul secțiunilor~\ref{sec:threat} și \ref{sec:consistenta}.

\subsection{Limitele modelului clasic al perturbărilor adversariale}
\label{subsec:lp}

În procesarea imaginilor, exemplele adversariale sunt definite frecvent prin raportare la o vecinătate a exemplului original. Sunt considerate admisibile acele variante pentru care distanța față de original, calculată de obicei cu ajutorul unei norme \(L_p\), nu depășește un prag stabilit. Aplicarea directă a acestui model traficului de rețea este însă problematică, deoarece distanța dintre doi vectori de caracteristici nu exprimă, singură, dacă transformarea corespunzătoare poate fi realizată asupra traficului.

\textbf{Caracteristicile unui flux sunt dependente.}
Mai multe dintre valorile care descriu un flux sunt calculate din aceleași informații de bază. De exemplu, dimensiunea medie a pachetelor depinde de volumul de date și de numărul de pachete, iar debitul depinde de volum și de durată. Modificarea separată a acestor caracteristici poate încălca relațiile dintre ele. Prin urmare, un vector poate respecta limita de distanță impusă și totuși să nu corespundă unui flux care ar putea fi observat în practică.

\textbf{Atacatorul nu controlează toate caracteristicile.}
Unele valori sunt determinate de comportamentul sistemului cu care acesta comunică sau sunt calculate de senzor pe baza întregului schimb de pachete. În această categorie pot intra volumul și numărul pachetelor transmise în sensul de retur, precum și anumite caracteristici agregate ale conexiunii. Includerea lor într-o vecinătate matematică nu înseamnă că atacatorul le poate modifica direct. Modelul perturbării trebuie, așadar, să distingă între caracteristicile asupra cărora atacatorul poate interveni și cele care rezultă din factori aflați în afara controlului său.

\textbf{Transformările realizabile sunt adesea direcționale.}
O vecinătate definită printr-o normă permite, în general, atât creșterea, cât și reducerea valorilor caracteristicilor. Operațiile aplicabile traficului nu au însă întotdeauna această simetrie. Adăugarea unor date de umplutură crește volumul transmis, iar introducerea unor întârzieri mărește durata fluxului. Reducerea acelorași valori poate necesita eliminarea unor date sau schimbarea succesiunii de acțiuni, ceea ce riscă să afecteze funcționalitatea atacului.

În consecință, admisibilitatea unei perturbări nu poate fi stabilită exclusiv prin distanța față de vectorul inițial. Ea trebuie definită prin transformări realizabile asupra traficului, limitate la acțiunile disponibile atacatorului și aplicate astfel încât relațiile dintre caracteristici și funcționalitatea atacului să fie păstrate.

\subsection{Fluxurile asupra cărora se face măsurarea}\label{subsec:eligibil}

Măsurarea evaziunii presupune stabilirea prealabilă a fluxurilor asupra cărora are sens să fie făcută. Se folosesc exclusiv fluxurile care îndeplinesc simultan două condiții: sunt în realitate atacuri și au fost semnalate ca malițioase de model pe traficul nemodificat, indiferent dacă modelul a nimerit sau nu categoria exactă.

Ambele condiții sunt necesare, din motive diferite.

Prima elimină fluxurile pe care modelul nu le detectează nici înainte de perturbare. Un asemenea flux trece nedetectat oricum, independent de orice acțiune a atacatorului. Dacă ar fi inclus în calcul, un eșec preexistent al modelului ar fi contabilizat drept succes al atacatorului, iar vulnerabilitatea măsurată ar fi sistematic exagerată. Se măsoară ce reușește atacatorul să obțină, nu ce ratase modelul dinainte.

A doua condiție renunță deliberat la cerința ca tipul de atac să fie identificat corect. Un flux malițios încadrat într-o categorie greșită declanșează totuși mecanismul de alarmă, deci este detectat în sensul operațional al termenului. Lucrarea distinge, așadar, între \emph{semnalarea} unui atac și \emph{acuratețea identificării subtipului} acestuia, iar evaluarea adversarială privește exclusiv prima dintre cele două. Distincția nu este una formală: cele două criterii diferă, în datele utilizate, cu aproape un sfert dintre fluxuri, iar alegerea între ele schimbă numitorul tuturor rezultatelor raportate ulterior.

\subsection{Cele trei rezultate posibile}\label{subsec:rezultate}

Pentru un flux inclus în măsurare, aplicarea unei perturbări poate duce la exact trei situații, care se exclud reciproc și acoperă toate posibilitățile. Fie modelul păstrează eticheta corectă, fie o schimbă în „trafic legitim”, fie o schimbă într-o altă categorie de atac. Interpretarea lor este rezumată în tabelul~\ref{tab:trei}.

\begin{table}[!ht]
    \caption{Cele trei rezultate posibile ale unei perturbări}
    \centering
    \begin{tabular}{|c|l|l|p{4.4cm}|}
        \hline
        & \textbf{Ce spune modelul} & \textbf{Alarmă?} & \textbf{Semnificație} \\
        \hline
        (a) & categoria corectă & da, corect & Perturbarea nu a avut efect \\
        \hline
        (b) & trafic legitim & \textbf{nu} & \textbf{Evaziune: atacul trece nedetectat} \\
        \hline
        (c) & altă categorie de atac & da, greșit & Eroare de clasificare, dar nu breșă de securitate \\
        \hline
    \end{tabular}
    \label{tab:trei}
\end{table}


\vspace{2.0cm}
\emph{Rata de evaziune} reprezintă proporția fluxurilor perturbate care sunt reclasificate drept trafic benign, raportată la numărul total de fluxuri incluse în evaluare. Aceasta măsoară succesul perturbării din perspectiva evitării detecției.

Clasificarea unui flux perturbat într-o altă categorie de atac nu este considerată evaziune, deoarece fluxul continuă să fie identificat drept malițios. Aceste confuzii rămân totuși relevante pentru analiza comportamentului clasificatorului. Frecvența și direcția lor arată ce clase sunt dificil de separat după aplicarea perturbării și permit identificarea tipurilor de atac între care modelul produce cel mai des confuzii.


\subsection{Proprietăți ale ratei de evaziune}\label{subsec:proprietati}

Interpretarea ratei de evaziune depinde de câteva aspecte metodologice relevante pentru comparațiile prezentate în capitolul~\ref{ch:testare}.

\textbf{Numitorul este specific fiecărui model.}
Setul de fluxuri inclus în calcul este stabilit separat pentru fiecare model și cuprinde atacurile pe care acesta le identifică drept malițioase înainte de perturbare. Prin urmare, ratele de evaziune ale unor modele diferite pot avea numitori diferiți. Această particularitate rezultă din definiția metricii: pot fi considerate evadate numai fluxurile detectate înainte de aplicarea perturbării.

Utilizarea unui set comun, format din fluxurile detectate de toate modelele comparate, ar facilita comparația directă, dar ar restrânge evaluarea fiecărui model la un subset influențat de performanța celorlalte. Din acest motiv, rata de evaziune este calculată separat pentru fiecare model, iar numărul fluxurilor incluse în calcul este raportat împreună cu valoarea obținută.

\textbf{Nivelurile de intensitate nu formează întotdeauna o scară crescătoare.} Pentru perturbările definite prin valori-țintă, cum este modificarea TTL, nivelurile succesive reprezintă valori distincte, nu intensități tot mai mari ale aceleiași modificări. Nu există, prin urmare, nicio garanție că rata de evaziune crește de la un nivel la altul: ea poate scădea, dacă valoarea-țintă a nivelului superior se întâmplă să fie mai apropiată de profilul traficului malițios decât cea a nivelului precedent. Interpretarea unei asemenea scăderi ca anomalie experimentală ar fi greșită. Pentru perturbările definite prin factori de multiplicare, cum sunt umplutura și dilatarea temporală, nivelurile sunt efectiv ordonate, iar creșterea este de așteptat.

\textbf{Rata de evaziune și rata de detecție sunt complementare.} Pe fluxurile incluse în măsurare, cazurile (a) și (c) înseamnă împreună că alarma se declanșează, iar cazul (b) că nu se declanșează. Proporția fluxurilor încă detectate este deci exact complementul ratei de evaziune. Această observație simplă este cea care permite, în secțiunea~\ref{sec:sensibilitate}, exprimarea importanței caracteristicilor în aceleași unități ca evaziunea. Fără ea, cele două mărimi ar putea fi puse în legătură doar calitativ.

\section{Elemente de protocol relevante}\label{sec:protocoale}

Perturbările analizate în lucrare nu sunt alese arbitrar dintre valorile numerice ale rezumatului de flux, ci corespund unor acțiuni concrete asupra pachetelor. Justificarea lor necesită câteva elemente ale protocoalelor din care provin mărimile respective.

\subsection{Câmpul TTL și mecanismul de decrementare}\label{subsec:ttl}

Antetul unui pachet IP conține câmpul \textit{Time-to-Live}, un număr întreg reprezentat pe opt biți. Rolul lui este pur operațional: fiecare ruter prin care trece pachetul îi scade valoarea cu o unitate, iar pachetul a cărui valoare ajunge la zero este eliminat din rețea și nu mai este transmis mai departe. Mecanismul există pentru a împiedica circulația la nesfârșit a pachetelor prinse într-o buclă de rutare.

Din acest mecanism decurg trei proprietăți esențiale pentru lucrare.

\textbf{Valoarea inițială este aleasă de emițător.} Ea nu este impusă de protocol, ci de configurația sistemului de operare al mașinii care trimite pachetul. Sistemele uzuale pornesc de la valori precum 64, 128 sau 255. Orice program poate stabili o altă valoare printr-o opțiune obișnuită a interfeței de programare a socketurilor, fără privilegii speciale și fără nicio modificare a conținutului transmis.

\textbf{Valoarea observată de senzor codifică distanța, nu intenția.} Ceea ce înregistrează senzorul nu este valoarea inițială, ci cea rămasă după traversarea rețelei, adică valoarea de plecare din care s-a scăzut numărul de rutere parcurse. Doi emițători configurați identic, dar aflați la distanțe diferite față de senzor, produc valori diferite. Doi emițători complet diferiți, aflați la aceeași distanță și pornind de la aceeași valoare, produc valori identice. Câmpul descrie, așadar, o proprietate a topologiei rețelei, nu una a traficului însuși.

\textbf{Modificarea nu afectează funcționalitatea atacului.} Atâta timp cât valoarea rămâne suficient de mare pentru a acoperi numărul de rutere până la țintă, pachetul ajunge la destinație cu conținutul neschimbat. Aceasta este proprietatea care face din modificarea TTL o perturbare cu cost practic nul pentru atacator.

Ultima proprietate impune totodată o limită inferioară. Sub un anumit prag, pachetul riscă să fie eliminat înainte de a ajunge la destinație, iar atacul devine nefuncțional. Perturbările analizate respectă această limită, altfel ar încălca însăși condiția de validitate pe care se sprijină întregul studiu.

\subsection{Direcționalitatea traficului}\label{subsec:directionalitate}

O conexiune de rețea are două sensuri, iar senzorul calculează mărimi separate pentru fiecare: câte pachete și câți octeți a trimis inițiatorul conexiunii, câte a trimis destinatarul, ce valori TTL s-au observat în fiecare sens.

Această separare are o consecință directă asupra modelului de amenințare. Atacatorul controlează stiva de rețea a propriei mașini, deci poate stabili câmpurile pachetelor pe care le emite. Nu are însă niciun mijloc de a determina ce conțin antetele pachetelor generate de victimă: valoarea TTL a răspunsurilor este stabilită de sistemul de operare al acesteia și redusă de ruterele aflate pe drumul de întoarcere.

Distincția este esențială pentru interpretarea rezultatelor din capitolul~\ref{ch:testare}, deoarece una dintre mărimile pe care modelele se sprijină cel mai mult provine tocmai din traficul de retur.

Există totuși o categorie intermediară. Dacă atacatorul își întârzie deliberat propriile pachete, conexiunea se întinde pe o durată mai mare, iar răspunsurile victimei se distribuie și ele pe acest interval mai lung, chiar dacă victima nu își schimbă în niciun fel comportamentul. Mărimile temporale asociate traficului de retur se modifică, prin urmare, ca efect indirect al unei acțiuni a atacatorului, fără ca acesta să le poată stabili direct. Tratarea lor ca fiind complet nemodificabile ar produce variante inconsistente, iar tratarea lor ca fiind direct controlabile ar acorda atacatorului o putere pe care nu o are.

\subsection{Starea conexiunii și identitatea atacului}\label{subsec:stare}

Senzorul înregistrează și starea în care s-a încheiat conexiunea, derivată din secvența indicatorilor din antetele TCP: o conexiune stabilită și închisă normal, una respinsă de destinatar, una rămasă deschisă sau una pentru care s-a observat doar o încercare de inițiere. Alături de protocolul folosit și de serviciul contactat, aceste valori descriu natura interacțiunii dintre cele două calculatoare.

Modificarea lor nu este considerată o perturbare admisibilă. Motivul nu este tehnic, ci logic: o încercare de conectare respinsă și o conexiune stabilită cu succes sunt evenimente diferite, nu variante ale aceluiași eveniment. Schimbarea unei asemenea valori nu ar produce o variantă perturbată a aceluiași atac, ci descrierea unui alt atac, iar compararea rezultatului cu eticheta reală ar deveni lipsită de sens.

\subsection{Segmentarea și limita fizică a umpluturii}\label{subsec:mtu}

Fiecare segment de rețea impune o dimensiune maximă a pachetului care poate fi transmis fără a fi fragmentat. Un atacator care adaugă octeți de umplutură nu poate, prin urmare, să mărească oricât pachetele existente: peste această limită este obligat să trimită pachete suplimentare.

Consecința asupra rezumatului de flux este că adăugarea de umplutură nu modifică o singură mărime. Peste un anumit prag, ea crește și numărul de pachete, iar odată cu acesta toate mărimile calculate din numărul de pachete. Este un al doilea exemplu, alături de cel din subsecțiunea~\ref{subsec:distanta}, în care o acțiune unică a atacatorului se propagă asupra mai multor mărimi simultan, și încă un motiv pentru care perturbarea trebuie construită la nivelul acțiunii, nu la nivelul valorii numerice.

\section{Modelul de amenințare}\label{sec:threat}

Modelul de amenințare precizează trei elemente: ce poate modifica atacatorul, ce știe despre sistemul de detecție și ce condiții trebuie să respecte modificările pentru a rămâne semnificative.

\subsection{Partiționarea caracteristicilor după controlabilitate}\label{subsec:taxonomie}

Punctul de plecare este observația că un atacator controlează propriile pachete, dar nu și pe cele emise de victimă. Această constatare, aparent evidentă, are consecințe cantitative importante, deoarece o parte din informația pe care modelul o folosește în decizia sa se află, prin construcție, în afara controlului atacatorului. Cele patruzeci și două de mărimi se împart în șapte categorii, prezentate în tabelul~\ref{tab:taxonomie}. Șase dintre ele intervin efectiv în construcția perturbărilor și sunt reprezentate schematic în figura~\ref{fig:taxonomie}. A șaptea cuprinde cele opt mărimi pe care niciuna dintre transformările analizate nu le atinge.

\begin{table}[!ht]
    \centering
    \begin{tabular}{|l|c|p{7.2cm}|}
        \hline
        \textbf{Categorie} & \textbf{Nr.} & \textbf{Ce reprezintă și cum este tratată} \\
        \hline
        Controlate direct de sursă & 7 & Calculate exclusiv din pachetele emise de atacator: TTL-ul trimis, volumul, numărul de pachete, durata. Modificabile în limitele fizice ale protocolului \\
        \hline
        Contoare de conexiuni & 7 & Numără câte dintre conexiunile recente au aceeași sursă, destinație sau serviciu ca fluxul curent. Descriu contextul, nu fluxul însuși. Pot fi doar reduse, prin încetinirea atacului \\
        \hline
        Derivate & 5 & Obținute din raportul altor mărimi, precum dimensiunea medie a pachetului. Nu se stabilesc direct, ci se recalculează prin propagare \\
        \hline
        Generate de victimă & 8 & Calculate din pachetele de răspuns: TTL-ul de retur, volumul și numărul de pachete primite. Constante, fiind în afara controlului atacatorului \\
        \hline
        Dependente de durată & 3 & Tot din răspunsurile victimei, dar se întind odată cu fluxul: se modifică numai ca efect al întârzierii introduse de atacator \\
        \hline
        Definitorii pentru atac & 4 & Protocolul, serviciul și starea conexiunii. Constante: modificarea lor ar însemna alt atac \\
        \hline
        Neperturbate & 8 & În afara domeniului transformărilor analizate \\
        \hline
    \end{tabular}
    \caption{Partiționarea caracteristicilor de flux după controlabilitatea lor de către atacator. Primele două categorii sunt singurele pe care atacatorul le poate modifica direct; următoarele două se recalculează automat, iar ultimele trei rămân neschimbate.}
    \label{tab:taxonomie}
\end{table}

\begin{figure}[!ht]
    \centering
    \begin{tikzpicture}[node distance=5mm and 8mm]
        \node[bloc, minimum width=32mm] (sursa) {Controlate direct\\ de sursă \;(7)};
        \node[bloc, minimum width=32mm, below=of sursa] (contoare) {Contoare de\\ conexiuni \;(7)};
        \node[bloc, minimum width=32mm, right=26mm of sursa] (derivate) {Derivate \;(5)};
        \node[bloc, minimum width=32mm, right=26mm of contoare] (cauzal) {Dependente de\\ durată \;(3)};
        \node[bloc, minimum width=32mm, right=22mm of derivate] (victima) {Generate de\\ victimă \;(8)};
        \node[bloc, minimum width=32mm, right=22mm of cauzal] (identitate) {Definitorii pentru\\ atac \;(4)};

        \draw[sageata] (sursa) -- node[eticheta, above] {propagare} (derivate);
        \draw[sageata] (contoare) -- node[eticheta, above] {propagare} (cauzal);

        \begin{scope}[on background layer]
            \node[grup, fit=(sursa)(contoare), label={[eticheta]above:\textbf{modificabile}}] {};
            \node[grup, fit=(derivate)(cauzal), label={[eticheta]above:\textbf{recalculate}}] {};
            \node[grup, fit=(victima)(identitate), label={[eticheta]above:\textbf{imuabile}}] {};
        \end{scope}
    \end{tikzpicture}
    \caption{Cele șase categorii care intervin în construcția perturbărilor, însumând 34 dintre cele 42 de caracteristici; restul de opt rămân în afara oricărei transformări și nu sunt reprezentate. Atacatorul acționează doar asupra grupului din stânga. Grupul din mijloc se modifică exclusiv prin propagare, iar cel din dreapta rămâne neschimbat.}
    \label{fig:taxonomie}
\end{figure}

Categoriile din tabelul~\ref{tab:taxonomie} se justifică după cum urmează.

\textbf{Caracteristicile controlate direct} sunt cele calculate exclusiv din pachetele emise de atacator: valoarea TTL a pachetelor trimise, volumul și numărul de pachete transmise, durata fluxului și parametrii ferestrei TCP. Modificarea lor corespunde unor acțiuni concrete la nivelul stivei de rețea a mașinii atacatoare.

\textbf{Contoarele de conexiuni} (denumite \textit{connection features} în documentația setului de date) numără câte dintre ultimele conexiuni înregistrate de senzor au aceeași sursă, aceeași destinație sau același serviciu ca fluxul analizat. Particularitatea lor este că nu descriu fluxul în sine, ci contextul în care acesta apare: două fluxuri identice în toate celelalte privințe primesc valori diferite dacă unul este izolat, iar celălalt face parte dintr-o serie de conexiuni către aceeași gazdă. Prefixul \verb|ct| din numele lor provine de la \textit{count}, iar sufixul \verb|ltm| de la \textit{last time}.

Ele sunt asimetrice din punctul de vedere al controlabilității: un atacator poate să le scadă, distribuind atacul în timp, dar nu le poate crește fără a genera trafic suplimentar, el însuși observabil. În consecință, perturbările asupra acestei categorii sunt admise doar în sens descrescător.

\textbf{Caracteristicile derivate} se calculează din altele și nu sunt niciodată stabilite direct de generator, ci recalculate în urma modificării mărimilor de bază, conform mecanismului din secțiunea~\ref{sec:consistenta}.

\textbf{Caracteristicile generate de victimă} (între care valoarea TTL a pachetelor de răspuns, volumul și numărul de pachete primite) nu pot fi stabilite de atacator. Ele constituie o rezervă de informație inaccesibilă, iar măsura în care un model se sprijină pe ea determină o limită naturală a evaziunii realizabile împotriva acelui model.

\textbf{Caracteristicile dependente de durată} reprezintă cazul intermediar discutat în subsecțiunea~\ref{subsec:directionalitate}. Ele sunt generate tot de victimă, dar nu sunt independente de acțiunea atacatorului: dacă acesta întârzie propriile pachete, răspunsurile victimei se întind și ele în timp, iar mărimile calculate din ele se modifică în consecință. Nu pot fi stabilite direct, ci doar prin propagare din raportul duratelor.

\textbf{Caracteristicile definitorii pentru atac} (protocolul, serviciul și starea conexiunii) sunt menținute constante, din motivele expuse în subsecțiunea~\ref{subsec:stare}.

\subsection{Accesul atacatorului la sistemul de detecție}\label{subsec:acces}

Se presupune că atacatorul nu dispune de parametrii modelului, de informația de gradient, de scorurile de încredere asociate predicțiilor și nici de o interfață prin care să poată interoga modelul, fie și numai pentru a obține eticheta prezisă. Perturbările sunt construite exclusiv pe baza proprietăților statistice observabile ale traficului, adică pe baza cunoașterii modului în care arată traficul legitim într-o rețea, nu pe baza vreunei reacții a sistemului de detecție.

Această ipoteză este cea mai restrictivă dintre variantele uzuale și, din acest motiv, produce estimări conservatoare: un atacator cu acces suplimentar ar obține rezultate cel puțin la fel de bune. Alegerea corespunde situației practice în care sistemul de detecție se află în infrastructura apărătorului și nu expune niciun canal de răspuns către exterior.

Ipoteza are însă și o consecință metodologică ce depășește realismul scenariului. Întrucât perturbarea nu depinde în niciun fel de model, exact aceleași variante perturbate pot fi aplicate tuturor modelelor evaluate. Comparația dintre modele devine astfel un experiment controlat: intrarea este identică, iar singura mărime care variază este modelul însuși. Orice diferență observată între ratele de evaziune este atribuibilă modelului, nu unei diferențe între perturbările la care a fost supus. Dacă perturbările ar fi fost optimizate separat pentru fiecare model, această atribuire nu ar mai fi fost posibilă.

Metodele care exploatează informația de gradient, precum FGSM \cite{goodfellow2015explaining}, se situează în afara acestui model de amenințare. Ele presupun accesul la derivatele funcției de decizie și, suplimentar, sunt aplicabile doar modelelor diferențiabile, ceea ce exclude ansamblurile de arbori. Chiar și atunci când sunt aplicabile, aceste metode produc perturbări direct în spațiul valorilor numerice, fără garanția că rezultatul corespunde unui trafic realizabil \cite{zhang2022adversarial, he2023adversarial}. Excluderea lor nu este, prin urmare, o afirmație despre inaplicabilitatea lor generală, ci o delimitare a cadrului adoptat aici.

\subsection{Mulțimea perturbărilor admisibile}\label{subsec:admisibile}

Reunind elementele anterioare, mulțimea perturbărilor admisibile se definește prin patru condiții simultane.

\begin{enumerate}
    \item \textbf{Direcționalitate.} Fiecare tip de transformare acționează într-un singur sens, determinat de acțiunea fizică pe care o reprezintă: volumul transmis poate doar să crească, durata poate doar să crească, contoarele de conexiuni pot doar să scadă.
    \item \textbf{Limite de domeniu relative.} Valorile rezultate trebuie să rămână în limitele fizic realizabile ale protocolului. Aceste limite se aplică relativ la valorile fluxului original, nu ca praguri absolute stabilite teoretic, deoarece capturile reale conțin ele însele înregistrări care încalcă limitele nominale. Impunerea unor praguri absolute ar respinge fluxuri care există efectiv în date.
    \item \textbf{Imuabilitate.} Caracteristicile generate de victimă și cele definitorii pentru identitatea atacului rămân neschimbate.
    \item \textbf{Închidere la dependențe.} Orice modificare a unei caracteristici de bază este însoțită de recalcularea tuturor caracteristicilor derivate din ea, astfel încât varianta rezultată să fie intern consistentă.
\end{enumerate}

Condiția a patra ridică cele mai multe dificultăți și face obiectul secțiunii următoare.

\section{Consistența internă a vectorului de caracteristici}\label{sec:consistenta}

\subsection{Problema inconsistenței induse}\label{subsec:inconsistenta}

Se consideră o perturbare care crește volumul de octeți transmiși prin adăugarea de umplutură în conținutul pachetelor. Dacă se modifică doar mărimea corespunzătoare, restul rezumatului rămâne cel calculat pentru volumul inițial. Rezultatul descrie un flux în care volumul transmis a crescut, dar dimensiunea medie a pachetului și debitul au rămas neschimbate, o combinație pe care niciun senzor de rețea nu ar putea să o producă, întrucât cele trei mărimi sunt legate prin relații deterministe.

Consecința asupra măsurătorii este severă. Dacă modelul își schimbă decizia pentru un asemenea flux, nu se poate stabili dacă a reacționat la perturbarea intenționată sau la imposibilitatea aritmetică a intrării. În al doilea caz, rata de evaziune măsurată nu ar caracteriza robustețea modelului la trafic modificat, ci comportamentul lui pe intrări situate în afara oricărei distribuții realiste. Consistența internă este deci o condiție de validitate a întregului experiment.

\subsection{Propagarea prin rapoarte}\label{subsec:propagare}

Soluția directă ar fi recalcularea caracteristicilor derivate din formulele documentate ale senzorului. Această abordare se lovește însă de o dificultate practică: formulele documentate nu reproduc exact valorile stocate în setul de date, întrucât constantele de conversie și convențiile de rotunjire folosite de senzor nu sunt pe deplin specificate. Recalcularea ar înlocui valoarea originală cu una ușor diferită, introducând o abatere în toate fluxurile, inclusiv în cele neperturbate, ceea ce ar încălca invariantul identității.

Se adoptă, în schimb, o abordare bazată pe rapoarte, justificată de următoarea observație.

\medskip
\noindent\textbf{Observația centrală.} \textit{Fie o caracteristică derivată care se obține împărțind o mărime de bază la alta și înmulțind rezultatul cu o constantă necunoscută, dar aceeași pentru toate fluxurile. Dacă perturbarea modifică mărimile de bază, atunci valoarea nouă a caracteristicii derivate se obține luând valoarea ei originală, înmulțind-o cu raportul dintre valoarea nouă și cea veche a mărimii de la numărător și împărțind-o la raportul corespunzător al mărimii de la numitor. Rezultatul este exact, indiferent de valoarea constantei necunoscute.}

\medskip

Constanta necunoscută intervine în același fel atât în valoarea originală, cât și în cea nouă. Atunci când valoarea nouă este exprimată prin cea veche, constanta apare o dată la numărător și o dată la numitor, deci se simplifică. Nu este nevoie, prin urmare, ca ea să fie determinată.

Această proprietate are două consecințe practice. Prima este că nu este necesară reconstituirea prin inginerie inversă a convențiilor senzorului: relațiile pot fi propagate corect fără a le cunoaște constantele. A doua, mai importantă, este că valoarea originală apare ca factor în rezultat. Pentru o perturbare nulă, ambele rapoarte devin egale cu unitatea, iar procedura redă exact valoarea stocată inițial, nu o aproximare a ei. Se păstrează astfel invariantul identității, care ar fi fost încălcat de recalcularea din formulă.

Modul în care se propagă modificările este prezentat în tabelul~\ref{tab:propagare}.

\begin{table}[!ht]
    \caption{Modul în care se propagă modificările asupra caracteristicilor derivate}
    \centering
    \begin{tabular}{|l|p{7.2cm}|}
        \hline
        \textbf{Caracteristică derivată} & \textbf{Cum se modifică} \\
        \hline
        Dimensiunea medie a pachetului & Crește proporțional cu volumul transmis și scade proporțional cu numărul de pachete \\
        \hline
        Debitul sursei & Crește proporțional cu volumul transmis și scade proporțional cu durata \\
        \hline
        Rata globală de pachete & Crește proporțional cu numărul total de pachete și scade proporțional cu durata \\
        \hline
        Intervalul mediu dintre pachete & Crește proporțional cu durata și scade proporțional cu numărul de intervale dintre pachete \\
        \hline
        Fluctuația temporală & Se modifică proporțional cu durata \\
        \hline
        Debitul destinației & Scade proporțional cu durata, volumul primit rămânând neschimbat \\
        \hline
    \end{tabular}
    \label{tab:propagare}
\end{table}

Cazurile degenerate necesită tratament explicit. Dacă o mărime de bază are valoarea zero în fluxul original, raportul corespunzător nu poate fi calculat, iar modificarea nu poate fi propagată. Situația apare pentru fluxurile formate dintr-un singur pachet, pentru care numărul de intervale dintre pachete este nul, și pentru cele cu durată înregistrată zero. În aceste cazuri caracteristica derivată se lasă neschimbată. Alegerea este conservatoare: ea reduce amplitudinea perturbării efective, deci nu poate produce o supraestimare a evaziunii.

\subsection{Caracteristici nereconstituibile și rezultate sub formă de interval}\label{subsec:interval}

Nu toate caracteristicile derivate se supun regulii de mai sus. Un caz distinct îl constituie contorul care combină starea conexiunii cu valorile TTL observate, numit în continuare \emph{contorul de context}, întrucât valoarea lui descrie situația în care apare fluxul, nu fluxul însuși. Analiza datelor arată că acest contor nu este o funcție a valorilor TTL ale fluxului curent: există perechi de fluxuri cu valori TTL apropiate cărora le corespund valori diferite ale contorului, precum și perechi cu valori TTL foarte diferite cărora le corespunde aceeași valoare. Comportamentul este consistent cu definiția sa reală, aceea de contor calculat pe o fereastră glisantă peste conexiunile recente.

Rezultă că mărimea nu poate fi calculată pornind de la un flux izolat: ea depinde de contextul în care fluxul apare, iar acest context nu este disponibil în reprezentarea tabelară. În consecință, valoarea pe care senzorul ar calcula-o pentru un flux perturbat nu poate fi determinată din datele existente.

Această nedeterminare nu poate fi eliminată, dar poate fi delimitată. Se definesc două politici explicite, corespunzătoare unor ipoteze extreme privind comportamentul contorului:

\begin{description}[style=nextline]
    \item[Politica conservatoare.] Contorul rămâne la valoarea din fluxul original. Ea corespunde ipotezei că modificarea TTL nu influențează contextul agregat, ceea ce subestimează efectul perturbării.
    \item[Politica optimistă.] Contorul adoptă valoarea tipică traficului legitim având aceeași valoare TTL. Ea corespunde ipotezei că atacatorul reușește să reproducă integral semnătura de context a traficului normal, ceea ce supraestimează efectul perturbării.
\end{description}

Evaluând rata de evaziune sub ambele politici se obține un interval, raportat ca atare în locul unei valori unice.

Natura acestui interval trebuie precizată, pentru a evita o interpretare greșită. El nu exprimă incertitudine statistică și nu are semnificația unui interval de încredere: nu provine din variabilitatea de eșantionare și nu i se poate asocia un nivel de acoperire. El reprezintă amplitudinea rezultatului între două ipoteze explicit formulate privind o mărime care nu poate fi reconstituită din date. Lățimea intervalului este ea însăși informativă: un interval îngust arată că rezultatul nu depinde practic de ipoteza adoptată, în timp ce un interval larg semnalează că o parte substanțială a efectului măsurat este condiționată de o presupunere care nu poate fi verificată.

Aceeași logică a delimitării explicite se aplică și caracteristicilor aflate în afara controlului atacatorului. Pe lângă perturbările realizabile, se poate evalua un scenariu de referință cu constrângeri relaxate, în care se modifică și caracteristicile generate de victimă. Un asemenea scenariu nu descrie un atac practic și nu poate fi raportat alături de rezultatele realizabile, rolul său este strict de a cuantifica ce parte din dependența unui model de o familie de caracteristici se sprijină pe informație inaccesibilă atacatorului.

\section{Algoritmii propuși}\label{sec:algoritmi}

Principiile prezentate anterior sunt puse în practică în patru proceduri, descrise independent de detaliile de implementare. Succesiunea etapelor urmărește respectarea constrângerilor definite în secțiunile precedente.

\subsection{Generarea unei variante perturbate}\label{subsec:alg-generare}

Procedura primește un flux și o transformare și produce fie o variantă validă, fie un refuz. Ordinea pașilor nu este arbitrară: propagarea trebuie să urmeze modificării de bază, iar validarea trebuie să fie ultima, altfel ar verifica o stare intermediară.

\begin{algorithm}[!ht]
\caption{Generarea unei variante perturbate}\label{alg:generare}
\begin{enumerate}\itemsep2pt
    \item Se pornește de la fluxul original și de la tipul și intensitatea transformării.
    \item Se stabilesc caracteristicile de bază vizate de transformare: valoarea TTL, volumul transmis, durata sau contoarele de conexiuni, respectând sensul unic admis pentru fiecare.
    \item Dacă transformarea a modificat volumul, se verifică limita de dimensiune a pachetului; la depășire, se adaugă pachete în loc să fie încălcată limita.
    \item Se recalculează toate caracteristicile derivate din cele modificate, prin procedura din subsecțiunea~\ref{subsec:alg-propagare}.
    \item Dacă valoarea TTL s-a schimbat, se stabilește și contorul de context. Acesta nu poate fi calculat din fluxul izolat, deci nu se recalculează, ci se alege conform uneia dintre cele două politici definite în subsecțiunea~\ref{subsec:interval}. Politica folosită se consemnează împreună cu varianta, întrucât de ea depinde marginea de interval pe care o va produce măsurarea.
    \item Se verifică cele șase condiții de validitate. Dacă vreuna este încălcată, varianta este \textbf{respinsă}, nu corectată.
    \item Se returnează varianta împreună cu evidența modificărilor efectuate.
\end{enumerate}
\end{algorithm}

Pasul 6 merită subliniat: o variantă care încalcă o constrângere nu este ajustată tacit până devine acceptabilă, ci semnalată ca eroare. Corectarea automată ar produce o variantă diferită de cea cerută, iar rezultatul măsurat nu ar mai corespunde parametrilor raportați.

\subsection{Propagarea caracteristicilor derivate}\label{subsec:alg-propagare}

Procedura pune în aplicare observația centrală din subsecțiunea~\ref{subsec:propagare}. Ea nu recalculează caracteristicile derivate din formulele senzorului, ci le modifică proporțional cu rapoartele dintre valorile noi și cele vechi ale mărimilor din care provin.

\begin{algorithm}[!ht]
\caption{Propagarea prin rapoarte}\label{alg:propagare}
\begin{enumerate}\itemsep2pt
    \item Pentru fiecare mărime de bază modificată se calculează raportul dintre valoarea nouă și cea originală.
    \item Dacă valoarea originală este nulă, raportul nu poate fi calculat; caracteristica derivată se lasă neschimbată.
    \item Fiecare caracteristică derivată se modifică proporțional cu rapoartele mărimilor din care se calculează, direct sau invers, după cum acestea apar la numărător sau la numitor.
    \item Mărimile temporale ale traficului de retur se modifică exclusiv cu raportul duratelor, conform discuției din subsecțiunea~\ref{subsec:directionalitate}.
    \item Se raportează numărul de fluxuri pentru care propagarea a fost neutralizată la pasul 2.
\end{enumerate}
\end{algorithm}

Pasul 5 permite identificarea situațiilor în care propagarea nu produce modificările așteptate. Dacă acest lucru se întâmplă pentru un număr mare de fluxuri, intensitatea efectivă a perturbării este mai redusă decât cea stabilită, iar efectul său poate fi subestimat.

\subsection{Măsurarea ratei de evaziune}\label{subsec:alg-masurare}

\begin{algorithm}[H]
\caption{Măsurarea ratei de evaziune}\label{alg:masurare}
\begin{enumerate}\itemsep2pt
    \item Se calculează predicțiile modelului pe traficul nemodificat și se determină fluxurile asupra cărora se face măsurarea, conform criteriului din subsecțiunea~\ref{subsec:eligibil}.
    \item Se verifică invariantul identității. Dacă rata obținută pe perturbarea nulă nu este exact zero, procedura se oprește: lanțul de măsurare este defect.
    \item Pentru fiecare transformare și fiecare nivel de intensitate se generează varianta corespunzătoare prin Algoritmul~\ref{alg:generare}.
    \item Se calculează predicțiile pe varianta perturbată, restrânse la fluxurile stabilite la pasul 1.
    \item Fiecare flux se încadrează în unul dintre cele trei rezultate din subsecțiunea~\ref{subsec:rezultate}, iar rata de evaziune se obține ca proporție a cazurilor de tipul (b).
    \item Pentru transformările cu două politici de context, cele două rate obținute delimitează intervalul discutat în subsecțiunea~\ref{subsec:interval}.
\end{enumerate}
\end{algorithm}

Pasul 2 este plasat înaintea evaluării, nu după ea, pentru că verifică însăși funcționarea procedurii. Perturbarea nulă nu modifică niciun flux, așa că o rată de evaziune nenulă în acest punct nu poate veni din perturbare, ci doar dintr-o defecțiune a lanțului de măsurare. Lăsată la final, aceeași verificare ar scoate defecțiunea la iveală abia după parcurgerea întregii grile, când rezultatele produse între timp ar fi consistente între ele și greu de deosebit de unele corecte.


\subsection{Analiza sensibilității}\label{subsec:alg-sensibilitate}

\begin{algorithm}[!ht]
\caption{Determinarea importanței și corelarea cu evaziunea}\label{alg:sensibilitate}
\begin{enumerate}\itemsep2pt
    \item Se măsoară proporția fluxurilor de atac pe care modelul le semnalează, pe date nemodificate.
    \item Pentru fiecare caracteristică, valorile ei se înlocuiesc cu valori extrase dintr-o distribuție exterioară, astfel încât legătura dintre caracteristică și etichetă să fie ruptă efectiv.
    \item Se remăsoară proporția de la pasul 1; scăderea observată constituie importanța caracteristicii. Substituția se repetă de mai multe ori, iar rezultatele se mediază.
    \item Se însumează importanțele pe categoriile de controlabilitate din tabelul~\ref{tab:taxonomie}, obținând proporția din capacitatea de detecție aflată la îndemâna atacatorului.
    \item Pentru fiecare transformare se însumează importanțele caracteristicilor pe care le afectează, iar seria rezultată se corelează cu ratele de evaziune măsurate.
\end{enumerate}
\end{algorithm}

Pasul 2 ascunde o subtilitate care s-a dovedit necesară în practică. Modul obișnuit de a rupe legătura dintre o caracteristică și etichetă este amestecarea valorilor ei în interiorul eșantionului analizat. Procedeul funcționează cât timp eșantionul este divers, dar se golește de conținut atunci când analiza se restrânge la o singură clasă de atac: dacă acolo caracteristica ia aproape aceeași valoare pentru toate observațiile, amestecarea le schimbă între ele fără să schimbe nimic, iar importanța rezultată apare nulă tocmai pentru o caracteristică pe care modelul se poate sprijini decisiv. Extragerea valorilor de înlocuire dintr-o distribuție exterioară eșantionului elimină acest artefact.

\section{Fundamentarea alegerii modelelor de clasificare}\label{sec:alegerea-modelelor}

Lucrarea compară trei modele: două ansambluri de arbori de decizie și o arhitectură neuronală bazată pe atenție. Alegerea nu urmărește acoperirea unui spectru cât mai larg de algoritmi, ci contrastarea a două familii cu moduri fundamental diferite de a construi decizia, întrucât tocmai această diferență este cea susceptibilă să se reflecte în comportamentul sub perturbări.

\subsection{Ansambluri de arbori de decizie}\label{subsec:arbori}

Ambele modele din prima familie construiesc decizia prin împărțiri succesive ale spațiului caracteristicilor, fiecare împărțire comparând o singură caracteristică cu un prag. Cele două sunt denumite în continuare prin numele lor consacrate, \textit{Random Forest} și \textit{XGBoost}. Denumirile descriptive, pădure aleatoare și creștere prin gradient, ar acoperi întreaga familie de algoritmi, nu implementările concrete ale căror rezultate sunt raportate aici.

\textbf{Random Forest} \cite{breiman2001random} combină predicțiile a sute de arbori, fiecare antrenat pe un eșantion diferit al datelor și fiecare alegându-și pragurile dintr-un subset de caracteristici extras aleator la fiecare nod. Cele două surse de aleatoriu îi fac pe arbori să greșească în locuri diferite, iar media predicțiilor anulează o parte dintre aceste erori fără să deplaseze decizia de ansamblu. Selecția aleatoare a caracteristicilor are și un efect care privește direct tema lucrării: fiindcă la fiecare nod doar o parte dintre caracteristici sunt disponibile, arborii nu ajung să se sprijine toți pe aceeași mărime dominantă, iar capacitatea de detecție se distribuie peste mai multe caracteristici.

\textbf{XGBoost} \cite{chen2016xgboost} construiește ansamblul secvențial, fiecare arbore nou fiind ajustat pe erorile rămase de la cei anteriori, prin optimizarea unei funcții obiectiv care include și un termen de regularizare. Spre deosebire de agregarea prin mediere, procedura reduce eroarea sistematică și tinde să exploateze intensiv caracteristicile cu putere predictivă ridicată. Din perspectiva teoretică a lucrării, aceasta conduce la o predicție verificabilă experimental: concentrarea capacității de detecție într-un număr mic de caracteristici ar trebui să crească vulnerabilitatea la o perturbare care vizează exact acele caracteristici.

Ambele modele au o proprietate comună cu consecințe metodologice directe: împărțirea spațiului folosește doar ordinea valorilor unei caracteristici, nu și magnitudinea acestora. Modelele sunt, prin urmare, insensibile la orice transformare care păstrează ordinea valorilor, iar aducerea numerelor la o scară comună nu le modifică deciziile. Aceeași proprietate explică însă și de ce o singură caracteristică poate domina decizia: un prag pe o caracteristică puternic corelată cu eticheta separă eficient clasele fără a necesita nicio combinație cu altele, iar mecanismul de construcție îl va prefera sistematic.

\subsection{Arhitectura bazată pe atenție pentru date tabelare}\label{subsec:atentie}

Arhitectura de tip transformer \cite{vaswani2017attention} a fost concepută pentru date secvențiale, precum textul, și nu se aplică direct datelor tabelare: nu există o secvență peste care să fie calculată atenția, iar ordinea coloanelor este arbitrară. Adaptarea folosită în lucrare urmează abordarea propusă de Gorishniy et al. \cite{gorishniy2021revisiting}, în care fiecare caracteristică este transformată într-un vector, iar mulțimea vectorilor astfel obținuți joacă rolul secvenței.

Concret, fiecare coloană numerică este transformată într-un vector printr-o operație liniară proprie, cu parametri învățați în timpul antrenării, iar fiecare coloană categorică primește un tabel de vectori indexat după valoarea categoriei. La secvența astfel obținută se adaugă un vector suplimentar, cu rol de agregator, din reprezentarea căruia se face în final clasificarea.

Peste această secvență se aplică blocuri succesive de auto-atenție. Mecanismul funcționează astfel: pentru fiecare pereche de elemente ale secvenței se calculează o măsură de asemănare, iar aceste măsuri sunt apoi normalizate, astfel încât ponderile asociate unui element să însumeze unitatea. În final, fiecare element este recalculat ca o combinație a tuturor celorlalte, ponderată cu valorile obținute. Operația se execută în paralel de mai multe ori, cu seturi diferite de parametri, ceea ce permite modelului să urmărească simultan mai multe tipuri de relații între caracteristici.

Diferența de mod de lucru față de ansamblurile de arbori este esențială pentru comparație. Mecanismul de atenție calculează ponderi pentru fiecare pereche de caracteristici, deci modelul reprezintă explicit relații între câte două mărimi, într-o formă continuă și învățată. Un ansamblu de arbori poate reprezenta asemenea relații doar implicit, prin înlănțuirea pragurilor pe ramuri succesive.

Această diferență atrage și o consecință metodologică. Spre deosebire de arbori, o rețea neuronală nu este insensibilă la scara valorilor: caracteristicile cu domenii de ordine de mărime diferite ar contribui disproporționat la ajustarea parametrilor. Modelul bazat pe atenție necesită, prin urmare, aducerea valorilor numerice la o scară comună, în timp ce ansamblurile de arbori nu o necesită. Este important de subliniat că această diferență de pregătire a datelor nu afectează validitatea comparației: mulțimea caracteristicilor, împărțirea în seturi de antrenare și testare și etichetele rămân identice pentru toate modelele. Aducerea la scară comună este o etapă internă unui singur model, nu o modificare a datelor comparate.

\subsection{Ce izolează comparația}\label{subsec:comparatie}

O comparație între modele spune ceva despre modele numai dacă tot restul rămâne neschimbat. Aici sunt ținute fixe patru lucruri: datele și împărțirea lor în seturi, mulțimea caracteristicilor, criteriul după care se aleg fluxurile pe care se măsoară și, pentru că perturbările nu depind în niciun fel de model, conform subsecțiunii~\ref{subsec:acces}, chiar variantele perturbate, identice rând cu rând pentru toate cele trei modele. Singura mărime care variază rămâne astfel modelul însuși, iar o diferență între ratele de evaziune se poate atribui felului în care fiecare își construiește decizia.

Comparația face însă și o distincție mai fină. Dacă un model cedează mai puțin decât celelalte, explicația poate fi una dintre două, iar ratele singure nu le deosebesc: fie modelul a învățat o reprezentare care nu se sprijină pe proprietăți ușor de modificat, fie se sprijină preponderent pe caracteristici la care atacatorul oricum nu ajunge. Prima ar fi o proprietate a arhitecturii, transferabilă la alte date, a doua, o consecință a felului în care au fost generate tocmai aceste date. Concluziile practice sunt opuse, iar separarea lor cere analiza de sensibilitate din secțiunea~\ref{sec:sensibilitate}.

\section{Evaluarea în condiții de dezechilibru sever}\label{sec:metrici}

Seturile de date pentru detecția intruziunilor sunt caracterizate de un dezechilibru pronunțat între clase \cite{ring2019survey, moustafa2015unsw}: categoriile de atac rare pot fi de câteva sute de ori mai puțin frecvente decât traficul normal. Alegerea metricilor de evaluare trebuie să țină cont de această situație.

Acuratețea globală, definită ca proporția predicțiilor corecte din total, este inadecvată. Într-un set în care o clasă acoperă o fracțiune majoritară din observații, un clasificator degenerat care prezice constant acea clasă obține o valoare ridicată fără a avea vreo capacitate de discriminare. Valoarea metricii este dominată de clasele frecvente, iar comportamentul pe clasele rare, chiar cele care corespund atacurilor mai puțin obișnuite devine practic invizibil.

Se folosesc, în consecință, metrici care acordă pondere egală fiecărei clase. Pentru fiecare categorie în parte se calculează trei mărimi. \emph{Precizia} arată ce proporție dintre fluxurile pe care modelul le-a atribuit acelei categorii îi aparțin în realitate, adică măsoară cât de multe alarme false produce modelul pentru categoria respectivă. \emph{Sensibilitatea} arată ce proporție dintre fluxurile care aparțin în realitate categoriei au fost efectiv recunoscute, adică măsoară cât de multe scapă nedetectate. \emph{Scorul F1} le combină pe amândouă într-o singură valoare, construită astfel încât să fie mică ori de câte ori una dintre cele două este mică, un model nu poate obține un scor bun sacrificând-o pe una în favoarea celeilalte.

Performanța globală se rezumă apoi prin media neponderată a scorurilor F1 ale tuturor categoriilor. Faptul că media este neponderată este esențial: el face ca eșecul pe o clasă rară să conteze exact la fel de mult ca eșecul pe una frecventă. Se raportează suplimentar media sensibilităților per clasă, precum și matricea de confuzie, care este singura reprezentare ce arată nu doar cât de des greșește modelul, ci și \emph{unde} plasează erorile, informație necesară pentru interpretarea rezultatului de tip (c).

Dezechilibrul este tratat și în etapa de antrenare, prin acordarea unei ponderi mai mari observațiilor din clasele rare, invers proporțională cu frecvența lor. În absența acestei corecții, contribuția claselor rare la mărimea pe care o optimizează procedura de antrenare devine neglijabilă, iar acestea ajung să fie ignorate.

O ultimă observație leagă această secțiune de criteriul din subsecțiunea~\ref{subsec:eligibil}. Un scor mediu moderat nu implică o capacitate slabă de semnalare a atacurilor. Cele două mărimi măsoară lucruri diferite: media penalizează încadrarea greșită a unui atac într-o categorie vecină, în timp ce semnalarea alarmei rămâne corectă în acel caz. Este perfect posibil ca un model să confunde frecvent categorii de atac apropiate ca profil, păstrând totodată o capacitate aproape completă de a distinge traficul malițios de cel legitim. Această observație justifică definirea mulțimii de măsurare pe criteriul semnalării, și nu pe cel al clasificării corecte: dacă cele două ar coincide, alegerea numitorului ar fi lipsită de consecințe.

\section{Fundamentarea analizei de sensibilitate}\label{sec:sensibilitate}

Măsurarea ratei de evaziune arată \emph{cât} cedează un model, dar nu explică \emph{de ce}. Explicația necesită o măsură a dependenței modelului de fiecare caracteristică, comparabilă între cele trei modele și exprimată în aceleași unități ca evaziunea.

\subsection{Necesitatea unei măsuri independente de model}\label{subsec:agnostic}

Fiecare dintre cele trei modele își raportează importanța altfel: Random Forest prin reducerea de impuritate, XGBoost prin câștigul acumulat în punctele de separare, iar modelul de tip Transformer prin ponderile mecanismului de atenție. Nu diferă doar formulele, ci și unitățile în care sunt exprimate, așa că valorile nu se pot pune una lângă alta. Cum aici se compară tocmai modelele între ele, e nevoie de o măsură unică, obținută în același fel pentru toate trei.

Importanța prin permutare este o astfel de măsură. Ea pornește de la o singură întrebare: ce se schimbă dacă o caracteristică nu mai spune nimic despre flux? Un model care se sprijinea pe ea ar trebui să înceapă să greșească vizibil, iar unul care nu o folosea ar trebui să nu simtă diferența. Se înlocuiesc, așadar, valorile acelei caracteristici cu valori străine de fluxul analizat, restul rezumatului rămâne neatins, iar cât pierde modelul din performanță devine chiar importanța caracteristicii. Nimic din procedură nu presupune ceva despre interiorul modelului și nimic nu cere reantrenare, deci cele trei modele înghețate pot fi trecute prin exact același test.

\subsection{Alegerea metricii și comparabilitatea}\label{subsec:comensurabil}

Alegerea mărimii a cărei degradare se măsoară este esențială. Se adoptă \emph{proporția fluxurilor malițioase încă semnalate ca atare}. Motivul este că această mărime este exact complementara celei măsurate prin rata de evaziune: evaziunea înseamnă că un flux de atac încetează să mai fie semnalat, iar scăderea proporției respective măsoară exact același fenomen.

În consecință, afirmațiile „modelul se sprijină în cutare măsură pe caracteristica X” și „perturbarea caracteristicii X produce o evaziune de cutare procent” sunt exprimate în aceleași unități și pot fi puse în corespondență direct, prin corelarea celor două serii de valori, nu doar prin analogie calitativă. Dacă mărimea aleasă ar fi fost, de exemplu, media scorurilor F1, cele două ar fi rămas legate doar intuitiv.

\subsection{Importanța accesibilă atacatorului}\label{subsec:accesibila}

Combinând măsura de mai sus cu partiționarea din tabelul~\ref{tab:taxonomie} se obține mărimea care leagă întreaga construcție teoretică. Se însumează importanțele tuturor caracteristicilor pe care atacatorul le poate modifica și se raportează rezultatul la suma importanțelor tuturor caracteristicilor. Se obține astfel o proporție, cuprinsă între zero și unu, numită în continuare \emph{fracțiunea accesibilă} a capacității de detecție. Importanțele negative, care apar atunci când amestecarea unei caracteristici îmbunătățește întâmplător rata de detecție, sunt trunchiate la zero înaintea însumării, atât la numărător, cât și la numitor. Fără această convenție, o sumă cu termeni de semne opuse ar putea produce proporții în afara intervalului sau instabile la variații mici ale măsurătorii.

Mărimea admite o interpretare directă: ea arată ce parte din informația pe care modelul o folosește efectiv pentru a semnala atacurile se află la îndemâna unui atacator. Un model cu o fracțiune accesibilă mică își sprijină decizia preponderent pe informație pe care atacatorul nu o poate atinge, ceea ce limitează evaziunea realizabilă independent de calitatea reprezentării învățate. Distincția introdusă în subsecțiunea~\ref{subsec:comparatie}, între rezistența datorată reprezentării și cea datorată inaccesibilității, devine astfel măsurabilă, iar nu doar enunțabilă.

\subsection{Limita atribuirii pe o singură caracteristică}\label{subsec:limita}

Procedura descrisă înlocuiește o singură caracteristică pe rând și, prin construcție, nu poate detecta efectele care apar doar la modificarea simultană și coerentă a mai multor caracteristici legate între ele. Dacă decizia modelului depinde de o combinație de valori, iar nu de fiecare valoare în parte, înlocuirea izolată a oricăreia dintre ele va produce o degradare modestă, sugerând în mod eronat o dependență redusă.

Această limitare are o consecință importantă pentru interpretarea rezultatelor: importanța atribuită individual constituie o limită inferioară a vulnerabilității adversariale, nu o estimare a ei. Ea justifică totodată, din nou, cerința de închidere la dependențe formulată în subsecțiunea~\ref{subsec:admisibile}: o metodologie care ar perturba caracteristicile izolat, fără propagare, ar rata tocmai efectele de acest tip.

\section{Antrenarea adversarială ca problemă de proiectare experimentală}\label{sec:antrenare}

Secțiunile anterioare au construit aparatul necesar pentru a măsura cât de ușor poate fi ocolit un sistem de detecție. Întrebarea complementară este dacă vulnerabilitatea astfel identificată poate fi redusă prin însuși modul în care modelul este antrenat. Ideea este simplă și veche: dacă modelul vede, în timpul antrenării, nu doar fluxurile de atac originale, ci și variante modificate ale acestora, etichetate în continuare drept malițioase, atunci este împins să recunoască atacul după proprietățile care rămân neschimbate sub transformare, nu după cele care se pot schimba.

Dificultatea nu stă în aplicarea ideii, ci în măsurarea corectă a efectului ei. Această secțiune arată de ce comparația directă dintre modelul reantrenat și cel inițial nu este suficientă și ce anume trebuie ținut fix pentru ca îmbunătățirea observată să poată fi atribuită perturbării.

\subsection{De ce augmentare și nu generare prin gradient}\label{subsec:de-ce-augmentare}

Există două moduri de a obține exemplele adăugate. În primul, ele sunt construite folosind derivatele funcției de decizie a modelului aflat în curs de antrenare, deci urmăresc granița de decizie pe măsură ce aceasta se deplasează. În al doilea, ele provin dintr-un set de transformări definite în prealabil, aplicate independent de reacția modelului.

Pentru cadrul construit aici, doar a doua variantă este disponibilă, din trei motive care se cumulează. Primul ține de simetria cu partea de atac: perturbările evaluate în subsecțiunea~\ref{subsec:lp} sunt construite fără acces la parametrii sau la răspunsurile modelului, iar o apărare antrenată pe exemple obținute prin gradient ar fi confruntată cu un atacator care nu dispune de aceeași informație. Comparația ar amesteca atunci două lucruri distincte, puterea metodei de generare și accesul presupus al atacatorului. Al doilea motiv ține de modelele analizate: ansamblurile de arbori nu admit derivate ale funcției de decizie în raport cu intrările, deci pentru două dintre cele trei clasificatoare metoda nu poate fi aplicată deloc, iar o apărare disponibilă doar pentru unul dintre ele nu ar permite compararea celor trei în aceleași condiții. Al treilea motiv, și cel mai important, este cel de realizabilitate: un exemplu obținut prin deplasarea vectorului de caracteristici în direcția indicată de derivate nu corespunde în mod necesar unui flux care poate fi generat efectiv, iar antrenarea pe astfel de exemple ar învăța modelul să reziste unor intrări care nu pot să apară în realitate.

Exemplele adăugate sunt, prin urmare, produse de același mecanism care generează variantele folosite la măsurarea evaziunii, cu aceleași verificări de admisibilitate și aceeași propagare a caracteristicilor derivate. Apărarea și atacul rămân astfel definite în raport cu aceeași mulțime de transformări, ceea ce face ca rezultatul să fie interpretabil: se măsoară rezistența la exact familia de modificări pentru care s-a demonstrat anterior vulnerabilitatea, nu la o familie mai largă și nici la una mai îngustă.

\subsection{Ce se schimbă simultan atunci când se adaugă exemple perturbate}\label{subsec:confuzii}

Copiile perturbate se adaugă fluxurilor originale, care rămân în setul de antrenare. Această alegere are o consecință care trebuie enunțată explicit, pentru că determină întreaga construcție experimentală: adăugarea unor copii doar pentru fluxurile de atac modifică simultan trei lucruri, nu unul singur.

Se modifică, în primul rând, cantitatea de exemple malițioase disponibile la antrenare. Un model antrenat pe mai multe exemple de atac poate detecta mai bine atacurile pentru simplul motiv că a văzut mai multe, independent de faptul că exemplele adăugate sunt sau nu perturbate. În al doilea rând, se modifică proporția dintre traficul normal și cel malițios. Într-un cadru dezechilibrat, această proporție influențează direct poziția granițelor de decizie, deci un efect apare chiar și în absența oricărei perturbări. În al treilea rând, dacă modelul folosește ponderi de clasă derivate din frecvențele observate, acestea se recalculează pe numărătorile modificate, deci se schimbă și funcția de cost minimizată, nu doar datele pe care este minimizată.

Un rezultat obținut prin compararea modelului reantrenat cu cel inițial nu poate distinge între aceste trei efecte și cel al perturbării propriu-zise. O scădere a ratei de evaziune ar fi observată și dacă exemplele adăugate ar fi simple copii ale celor existente, iar atribuirea ei perturbării ar fi nejustificată.

\subsection{Grupul de control și ce anume izolează}\label{subsec:control}

Soluția constă în adăugarea unui al treilea model, antrenat pe un set construit astfel încât să difere de cel al modelului apărat printr-o singură proprietate. Acest grup de control primește exact același număr de exemple de atac suplimentare, provenite din exact aceleași fluxuri sursă, dar neperturbate: simple duplicate ale înregistrărilor originale.

Prin construcție, cele două seturi au același număr de rânduri, aceeași distribuție a claselor și același efort de antrenare, măsurat în număr de exemple parcurse. Singura diferență între ele este dacă rândurile adăugate au fost sau nu modificate. Cele trei modele formează atunci o succesiune în care fiecare diferență are o interpretare precisă: distanța dintre modelul de referință și cel de control măsoară efectul volumului suplimentar de date și al deplasării proporției dintre clase, iar distanța dintre modelul de control și cel apărat măsoară efectul perturbării, care este mărimea de interes. Figura~\ref{fig:brate} prezintă această construcție.

Această construcție are și un cost pe care merită să îl anticipăm: ea triplează efortul de antrenare, întrucât fiecare configurație trebuie antrenată de trei ori, o dată pentru fiecare grup. Costul este însă condiția pentru ca afirmația finală să fie una cauzală și nu doar o observație de corelație.

\begin{figure}[!h]
    \centering
    \begin{tikzpicture}[node distance=7mm and 15mm]
        \node[blocmare] (s1) {Set original};
        \node[blocmare, below=of s1] (s2) {Set original\\ + duplicate neperturbate};
        \node[blocmare, below=of s2] (s3) {Set original\\ + copii perturbate};

        \node[bloc, right=of s1, minimum width=24mm] (m1) {Model de\\ referință};
        \node[bloc, right=of s2, minimum width=24mm] (m2) {Model de\\ control};
        \node[bloc, right=of s3, minimum width=24mm] (m3) {Model\\ apărat};

        \node[blocmare, right=17mm of m2] (masura) {Măsurare pe\\ mulțimea comună};

        \draw[sageata] (s1) -- (m1);
        \draw[sageata] (s2) -- (m2);
        \draw[sageata] (s3) -- (m3);
        \draw[sageata] (m1) -- (masura);
        \draw[sageata] (m2) -- (masura);
        \draw[sageata] (m3) -- (masura);

        \draw[punctat] (m1) -- node[eticheta, left] {efectul volumului} (m2);
        \draw[punctat] (m2) -- node[eticheta, left] {efectul perturbării} (m3);
    \end{tikzpicture}
    \caption{Cele trei grupuri ale experimentului. Seturile de antrenare ale ultimelor două au același număr de rânduri și aceeași distribuție a claselor, diferă doar prin faptul că rândurile adăugate sunt sau nu perturbate. Rezultatul raportat este diferența dintre modelul apărat și cel de control, nu față de modelul de referință.}
    \label{fig:brate}
\end{figure}

\subsection{Menținerea neschimbată a funcției de cost}\label{subsec:ponderi}

Al treilea efect enumerat în subsecțiunea~\ref{subsec:confuzii} nu este eliminat de grupul de control, pentru că se manifestă identic în ambele grupuri augmentate, dar el trebuie totuși neutralizat pentru ca acestea să rămână comparabile cu modelul de referință.

Modelele antrenate în condiții de dezechilibru sever folosesc ponderi de clasă invers proporționale cu frecvența observată, conform argumentului din secțiunea~\ref{sec:metrici}. Dacă aceste ponderi ar fi recalculate pe setul augmentat, ele ar avea alte valori decât cele folosite de modelul de referință, iar cele două modele ar diferi atât prin datele de antrenare, cât și prin funcția minimizată. Soluția adoptată este calcularea ponderilor o singură dată, pe etichetele setului original, și transmiterea aceluiași set de valori tuturor grupurilor. Augmentarea schimbă astfel numai datele, nu și obiectivul.


\subsection{Proveniența statisticilor folosite la generare}\label{subsec:provenienta}

Generatorul de variante are nevoie, pentru caracteristicile care nu pot fi recalculate din fluxul izolat, de statistici extrase din datele observate, conform delimitării din subsecțiunea~\ref{subsec:interval}. Atunci când modelele sunt înghețate, aceste statistici pot fi construite din ambele partiții: ele descriu comportamentul senzorului, adică o proprietate a procesului de generare a datelor, iar nimic nu se antrenează pe ele.

În momentul în care variantele perturbate devin date de antrenare, situația se schimbă calitativ. O statistică derivată din partiția de test s-ar regăsi atunci, indirect, în valorile pe care modelul le vede la antrenare, ceea ce ar constitui o scurgere de informație dinspre setul de evaluare către model. Din acest motiv, generarea datelor de antrenare folosește exclusiv statistici construite din partiția de antrenare.

Distincția este între două roluri diferite ale aceluiași mecanism, nu o inconsecvență. Generarea variantelor pe care se măsoară evaziunea continuă să folosească statisticile complete, pentru ca modificările măsurate să fie identice cu cele din etapele anterioare și rezultatele să rămână comparabile. Acolo, modelele sunt deja antrenate și nu mai pot fi influențate. Restricția se aplică numai generării datelor de antrenare. Consecința practică este că generatorul dispune, în acest al doilea rol, de un tabel de corespondențe mai sărac, deci recurge mai des la valorile de rezervă, iar acest lucru trebuie raportat.

\subsection{Numitorul comun al comparației}\label{subsec:cohorta}

Rata de evaziune se calculează, conform subsecțiunii~\ref{subsec:eligibil}, exclusiv peste fluxurile pe care modelul le semnalează pe traficul nemodificat. Această restricție este necesară, dar creează o dificultate atunci când se compară modele diferite: fiecare model semnalează o mulțime proprie de fluxuri, deci fiecare rată se raportează la o populație proprie.

Consecința este că două procente obținute de modele diferite nu sunt direct comparabile. Un model care semnalează mai puține fluxuri pe traficul nemodificat are un numitor mai mic. Dacă fluxurile pe care nu le mai semnalează sunt tocmai cele ușor de perturbat, rata lui de evaziune scade fără ca robustețea să fi crescut. Comparația ar înregistra atunci o îmbunătățire produsă de o pierdere de sensibilitate.

Soluția este fixarea în prealabil a unei mulțimi comune de fluxuri, formată din intersecția fluxurilor semnalate de toate modelele comparate, la toate repetările experimentului. Pe această mulțime, care nu se modifică de la un model la altul, ratele se măsoară pe exact aceleași fluxuri și devin comparabile pereche cu pereche. Mărimea mulțimii trebuie raportată împreună cu ratele: un procent calculat pe o mulțime restrânsă nu are aceeași greutate ca unul calculat pe o mulțime largă, iar o mulțime care s-a micșorat drastic este ea însăși un rezultat, semnalând că modelele comparate au ajuns să semnaleze fluxuri diferite.

Rămâne de precizat ce valoare se reține atunci când se compară două modele pe întreaga grilă de transformări. Media peste variante ar fi înșelătoare, întrucât atacatorul nu alege o transformare la întâmplare, ci pe cea care îi convine. Valoarea reținută este, prin urmare, maximul peste variantele admisibile: o apărare care reduce media, dar lasă neatinsă o singură variantă eficientă, nu a apărat nimic.

\subsection{Criterii stabilite înaintea rulării}\label{subsec:preinregistrare}

Ultimul element al construcției privește modul în care se decide dacă rezultatul este sau nu un succes. Dacă pragurile sunt alese după ce cifrele sunt cunoscute, ele pot fi ajustate, conștient sau nu, astfel încât rezultatul obținut să le satisfacă. Pragurile sunt de aceea fixate înainte de rulare și consemnate împreună cu restul configurației.

Trei condiții sunt formulate în acest fel. Prima privește câștigul de robustețe și este exprimată ca îmbunătățire față de grupul de control, nu ca prag absolut. Motivul este că un prag absolut ar putea marca drept eșec un rezultat valoros: o reducere a evaziunii de la o valoare foarte mare la una încă ridicată rămâne o observație importantă, chiar dacă nu atinge un nivel considerat acceptabil în practică. A doua condiție privește costul pe traficul nemodificat și limitează scăderea admisă a performanței globale. A treia este o condiție de siguranță și cere ca nicio clasă de atac să nu piardă mai mult decât un prag stabilit din capacitatea de a fi detectată. Ea există pentru că o apărare poate consolida clasa cea mai vizibilă degradând clasele rare, iar o metrică globală ar ascunde acest lucru.

Pe lângă aceste criterii este stabilită și o valoare de referință, care indică nivelul de performanță urmărit în practică, fără a condiționa validitatea experimentului. Dacă rezultatele nu ating criteriile stabilite, acest lucru arată că metoda de antrenare evaluată nu reduce suficient efectul transformărilor analizate. Un asemenea rezultat rămâne relevant, chiar dacă nu confirmă eficiența metodei.

\section{Structura logică și funcțională a soluției}\label{sec:structura}

Din elementele teoretice anterioare rezultă o descompunere funcțională în șase module, cu dependențe strict secvențiale. Fluxul de date este prezentat în figura~\ref{fig:pipeline}.

\begin{figure}[!ht]
    \centering
    \begin{tikzpicture}[node distance=8mm and 21mm]
        \node[bloc, minimum width=20mm] (date) {Date\\ UNSW-NB15};
        \node[blocmare, right=of date] (m1) {\textbf{M1}\\ Antrenare și evaluare};
        \node[blocmare, right=of m1] (m2) {\textbf{M2}\\ Generarea variantelor};
        \node[blocmare, below=13mm of m2] (m3) {\textbf{M3}\\ Măsurarea evaziunii};
        \node[blocmare, left=of m3] (m4) {\textbf{M4}\\ Analiza sensibilității};
        \node[bloc, left=of m4, minimum width=20mm] (m6) {\textbf{M6}\\ Interfața};
        \node[blocmare, below=13mm of m4] (m5) {\textbf{M5}\\ Antrenarea adversarială};

        \draw[sageata] (date) -- (m1);
        \draw[sageata] (m1) -- node[eticheta, above] {fluxurile de} node[eticheta, below] {măsurare} (m2);
        \draw[sageata] (m2) -- node[eticheta, right] {28 de variante} (m3);
        \draw[sageata] (m3) -- node[eticheta, midway, align=center, fill=white, inner sep=1pt] {rezultate\\ per flux} (m4);
        \draw[sageata] (m4) -- (m6);
        \draw[sageata] (m5.west) -| node[eticheta, pos=0.32, above] {rezultatele} node[eticheta, pos=0.32, below] {antrenării} (m6.south);
        \draw[sageata] (m2.east) -- ++(10mm,0) |- node[eticheta, pos=0.82, above] {generatoarele} (m5.south);
        \draw[sageata] (m5.east) -| node[eticheta, pos=0.28, above] {trei grupuri} (m3.south);
        \draw[sageata] (m5.north) -- node[eticheta, right] {modele apărate} (m4.south);

        \draw[punctat] (m1.south) -- node[eticheta, pos=0.62, left] {modele înghețate} (m4.north);
    \end{tikzpicture}
    \caption{Structura funcțională a soluției și fluxul de date între module. Săgeata punctată indică modelele înghețate din M1, folosite fără reantrenare de M4. M4 se aplică de două ori, o dată modelelor de referință și o dată celor apărate produse de M5. Interfața, M6, este singurul modul terminal, consumând atât rezultatele analizei de sensibilitate, cât și pe cele ale antrenării adversariale.}
    \label{fig:pipeline}
\end{figure}

\textbf{Modulul M1} construiește clasificatoarele și le evaluează pe traficul nemodificat, folosind metricile din secțiunea~\ref{sec:metrici}. El produce modelele antrenate și predicțiile de referință, din care se derivă mulțimea fluxurilor asupra cărora se va măsura evaziunea. După această etapă, modelele sunt considerate înghețate: niciun modul ulterior nu le modifică, ceea ce garantează că variațiile observate provin exclusiv din perturbarea intrărilor.

\textbf{Modulul M2} generează variantele perturbate. El pune în aplicare cele patru condiții din subsecțiunea~\ref{subsec:admisibile} și procedura de propagare din subsecțiunea~\ref{subsec:propagare}. Fiecare variantă generată este supusă unui set de verificări de validitate înainte de a fi transmisă mai departe. O variantă care încalcă o constrângere este respinsă, nu corectată tacit. Modulul depinde de M1 doar prin mulțimea fluxurilor de măsurat, nu prin structura modelelor.

\textbf{Modulul M3} aplică modelele înghețate variantelor perturbate și încadrează rezultatele conform taxonomiei din subsecțiunea~\ref{subsec:rezultate}, calculând ratele și intervalele. Verificarea invariantului identității precede orice altă măsurătoare. Mecanismul este invocat de două ori în cursul experimentului, o dată pentru modelele din M1 și o dată pentru cele din M5, iar cele două rulări scriu rezultate în locații separate, pentru ca ratele modelelor de referință să nu se amestece cu cele ale modelelor apărate.

\textbf{Modulul M4} calculează importanțele și fracțiunile accesibile și pune în coresponden\-ță importanța caracteristicilor afectate de fiecare transformare cu rata de evaziune produsă de acea transformare. El consumă atât modelele din M1, cât și rezultatele per flux din M3. Aceeași procedură se aplică, într-o a doua rulare, modelelor produse de M5, pentru a stabili dacă apărarea a schimbat caracteristicile pe care acestea se sprijină. Rezultatul acestei a doua rulări este prezentat în subsecțiunea~\ref{subsec:o6-mecanism}.

\textbf{Modulul M5} reia antrenarea clasificatoarelor pe seturi augmentate cu variante perturbate, conform construcției din secțiunea~\ref{sec:antrenare}. Spre deosebire de celelalte module, el nu consumă modelele înghețate, ci produce modele noi, motiv pentru care artefactele sale sunt păstrate separat și nu înlocuiesc niciodată modelele de referință. Dependența lui de M2 este una de mecanism, nu de rezultat: folosește aceleași generatoare, dar aplicate partiției de antrenare și cu statisticile restrânse conform subsecțiunii~\ref{subsec:provenienta}. Legătura cu M3 este, la rândul ei, tot una de mecanism, conform explicației de mai sus. Bucla astfel închisă nu introduce circularitate, întrucât măsurarea se face cu modelele deja antrenate și înghețate la rândul lor.

\textbf{Modulul M6} expune rezultatele într-o formă explorabilă interactiv și permite evaluarea unor combinații de parametri care nu fac parte din setul predefinit. Cerința teoretică impusă acestui modul este ca perturbările pe care le aplică să fie construite prin aceleași mecanisme ca cele din M2, iar nu printr-o realizare paralelă. Altfel, valorile afișate interactiv ar putea diverge de cele măsurate sistematic, iar modulul ar înceta să fie un instrument de explorare a acelorași rezultate. Fiind ultimul din lanț, este și singurul care afișează simultan rezultatele lui M4 și pe cele ale lui M5, motiv pentru care depinde de amândouă, nu doar de modulul care îl precede direct.

Dependențele dintre module sunt unidirecționale, iar fiecare tranziție este însoțită de o condiție verificabilă: varianta neperturbată generată de M2 trebuie să reproducă predicțiile de referință ale lui M1, iar rata de evaziune calculată de M3 pentru această variantă trebuie să fie nulă. Aceste condiții transformă lanțul de prelucrare într-o succesiune de etape verificabile independent, ceea ce permite localizarea unei eventuale erori la nivelul modulului care o produce.

\section{Ipotezele cadrului și limitele lor}\label{sec:ipoteze}

Cadrul teoretic construit se sprijină pe un număr de ipoteze care nu pot fi verificate în interiorul lui. Enunțarea lor explicită este necesară pentru delimitarea corectă a domeniului de valabilitate a rezultatelor.

\textbf{Ipoteza de realizabilitate.} Se presupune că orice variantă considerată admisibilă corespunde unui trafic pe care atacatorul îl poate genera efectiv. Condițiile din subsecțiunea~\ref{subsec:admisibile} sunt construite pentru a garanta acest lucru, dar garanția rămâne una deductivă: ea decurge din interpretarea fiecărei caracteristici, nu dintr-o verificare experimentală. O confirmare completă ar presupune generarea traficului modificat, trecerea lui printr-un senzor real și compararea rezumatului rezultat cu cel obținut prin propagare. Această verificare depășește cadrul lucrării, iar rezultatele trebuie citite ca fiind condiționate de ipoteză.

\textbf{Ipoteza privind forma relațiilor derivate.} Observația centrală din subsecțiun\-ea~\ref{subsec:propagare} presupune că fiecare caracteristică derivată se obține ca raport a două mărimi de bază, înmulțit cu o constantă necunoscută, dar aceeași pentru toate fluxurile. Ipoteza este verificabilă indirect, prin măsurarea degradării relațiilor dintre mărimi după propagare, dar nu poate fi demonstrată în absența specificației complete a senzorului. Acolo unde forma presupusă ar fi greșită, propagarea ar introduce o abatere sistematică, mărimea acestei abateri este însă observabilă și raportabilă.

\textbf{Ipoteza de invarianță a procedurii de extragere.} Se presupune că senzorul calculează caracteristicile în același mod pentru traficul original și pentru cel perturbat. Ipoteza este rezonabilă pentru mărimile calculate din fluxul însuși, dar devine problematică pentru cele care depind de context, motiv pentru care acestea din urmă sunt tratate prin delimitarea din subsecțiunea~\ref{subsec:interval} în loc de a fi recalculate.

\textbf{Ipoteza de reprezentativitate a setului de date.} Concluziile privind ce caracteristici sunt importante pentru un model se referă la modelul antrenat pe acest set, nu la clasa de modele în general. Dacă o caracteristică are, în setul utilizat, o corelație cu eticheta care provine din modul de generare a datelor și nu din natura atacurilor, atunci dependența modelului de acea caracteristică este un artefact al setului. Cadrul propus nu poate distinge singur între cele două situații, iar distincția necesită analiza separată a provenienței datelor.

\textbf{Ipoteza de neadaptivitate.} Modelul de amenințare din subsecțiunea~\ref{subsec:acces} exclude orice interacțiune între atacator și sistemul de detecție. Un atacator cu acces la răspunsurile modelului ar putea căuta perturbări mai eficiente decât cele dintr-un set predefinit. Rezultatele obținute în acest cadru constituie, prin urmare, estimări conservatoare ale vulnerabilității: ele indică ce se poate obține fără nicio cunoaștere a apărării, nu limita superioară a ceea ce se poate obține.

\textbf{Adaptarea atacatorului nu este evaluată.}
După reantrenare, modelul este testat folosind aceleași tipuri de transformări considerate anterior. Evaluarea nu analizează situația în care atacatorul observă efectul apărării și își modifică strategia. Prin urmare, reducerea ratei de evaziune arată că modelul a devenit mai rezistent la transformările evaluate, dar nu demonstrează robustețea sa față de alte strategii de atac.

\textbf{Generalizarea la transformări noi nu este evaluată separat.}
Utilizarea acelorași tipuri de transformări atât la antrenare, cât și la testare nu permite stabilirea măsurii în care modelul poate rezista unor perturbări diferite. O asemenea verificare ar necesita excluderea unui tip de transformare din setul de antrenare și utilizarea sa numai în etapa de evaluare. În absența acestui experiment, concluziile privind eficiența antrenării se limitează la familiile de transformări incluse în studiu.

Ultima observație privind neadaptivitatea este cea mai importantă pentru interpretarea corectă a rezultatelor. O rată de evaziune ridicată măsurată în acest cadru constituie o dovadă solidă de vulnerabilitate, întrucât a fost obținută în condițiile cele mai defavorabile atacatorului. O rată scăzută, în schimb, nu constituie o dovadă de robustețe, ci doar absența unei vulnerabilități la clasa specifică de transformări evaluate.
